\section{Introduction}\label{sec:introduction}

Let $M$ be a closed smooth surface: a compact smooth two-dimensional manifold
without boundary. We do not assume that $M$ is orientable or connected.
For a smooth positive definite metric $g$ on $M$, an isometric immersion into
Euclidean four-space is a smooth map $F:M\to\R^4$ satisfying
\[
 \langle dF_p(v),dF_p(w)\rangle=g_p(v,w)
 \qquad(p\in M,\ v,w\in T_pM).
\]
Positive definiteness then makes $dF_p$ injective at every point.
The question is whether every such metric can be realized in this dimension.

\begin{theorem}\label{thm:main}
Every closed smooth Riemannian surface $(M,g)$ admits a $C^\infty$ isometric
immersion into $(\R^4,\delta_4)$, where $\delta_4=\sum_{i=1}^4 dx_i^2$.
\end{theorem}

The result concerns immersion; self-intersections are allowed. There is no
restriction on Gaussian curvature or on the topology of the closed surface.
In particular, the construction applies directly to nonorientable surfaces.

\subsection{Context}
Gromov explicitly recorded the question of smooth isometric immersion of
arbitrary smooth Riemannian surfaces into $\R^4$
\cite[\S I, Remark~(d), p.~3 of the author version]{Gromov2000}.
He later attributed the four-space surface question to Chern around 1950,
while expressing uncertainty about the compactness qualification
\cite[p.~4, footnote~6]{Gromov2015}.
Theorem~\ref{thm:main} concerns the compact boundaryless case of this question.

Nash's smooth embedding theorem gives a global realization of a compact
$n$-dimensional Riemannian manifold in dimension $n(3n+11)/2$, which is $17$
for surfaces \cite[Theorem~2, p.~59]{Nash1956}.
A closer surface comparison is Gromov's construction of smooth
isometric immersions of closed surfaces into $\R^5$: his retrospective
explicitly recalls this result and refers to \S3.2.4 of
\emph{Partial Differential Relations}
\cite[p.~59]{Gromov2015}\cite{GromovPDR1986}.
Theorem~\ref{thm:main} lowers this immersion dimension to four.

The local analytic theory reaches a smaller dimension. The theorem of
Janet and Cartan places an analytic Riemannian surface locally and
isometrically in $\R^3$ \cite{Janet1926,Cartan1927}; see also the
formulation in \cite[\S1.2]{Gromov2015}.
Both the analytic hypothesis and the local conclusion matter here.
For comparison, the companion \cite[Theorem~A]{OpenAILocalObstruction2026}
constructs a smooth positive definite metric on $(-1,1)^2$ for which no
neighborhood of the origin admits a smooth isometric immersion into $\R^3$.

There is also a smooth four-dimensional theory on planar domains.
Poznyak proved that $C^{3,\alpha}$ data admit $C^{2,\alpha}$ isometric
immersions into $\R^4$ on compact portions of a simply connected planar
domain with regular boundary \cite[Theorem~4, p.~67]{Poznyak1973}.
The smooth version of this construction is stated and explained in
Lewicka's modern account \cite[Theorem~9.1, p.~57]{Lewicka2026v3}.
Passing from planar domains to a closed surface requires the corrections
to remain compatible through a global construction. Here we achieve this
by preserving the normal geometry needed for successive local metric
corrections.

Regularity gives a second distinction. An immersion is \emph{strictly short}
if the prescribed metric minus its induced metric is positive definite.
Nash's oscillatory construction and Kuiper's codimension-one refinement show
that a strictly short immersion in positive codimension can be uniformly
approximated by $C^1$ isometric
immersions \cite{Nash1954,Kuiper1955}; the contemporaneous account
\cite[Theorems~1--2]{Choquet1958} states the immersion form explicitly.
Lewicka proves a finer flexibility theorem in $\R^4$ for metrics of low
H\"older regularity. In the August 2026 version, the domain is bounded,
planar, and diffeomorphic to a disk, the metric is $C^{r,\beta}$ on its
closure with $0<r+\beta\le2$, and strictly short $C^1$ immersions admit
uniform approximation by $C^{1,\alpha}$ isometric immersions for
$0<\alpha<\min\{(r+\beta)/2,1\}$
\cite[Theorem~1.1, p.~2]{Lewicka2026v3}.
These domain and regularity hypotheses differ from the smooth
closed-surface statement of Theorem~\ref{thm:main}.

The proof uses the established ideas of positive rank-one metric
decomposition, oscillatory metric correction, and smoothing to control
derivative loss. We give the corresponding attribution where these ideas
enter the construction. The mode equations, finite correctors, smooth
convergence argument, and propagation of the normal geometry are proved
below. Whitney's ordinary immersion theorem supplies the starting map
\cite[Theorem~8, p.~274]{Whitney1944}; Sard's equal-dimensional
critical-value theorem is used to choose the support boundaries and phases
\cite[Theorem~4.1, p.~885]{Sard1942}.

\subsection{The construction}
The proof has four stages. First prepare a smooth immersion into a small round
three-sphere in $\R^4$ together with a finite decomposition of its metric
deficit. Its induced metric $g_0$ is strictly below $g$, and the negative
radial direction is a globally defined normal, including when $M$ is
nonorientable. The decomposition has the form
\begin{equation}\label{eq:intro-sum}
 g-g_0=\sum_{j=1}^{J_0} a_j^2\,dx_j^2,
\end{equation}
where each $x_j$ is a coordinate function on a disk and $a_j$ is smooth,
positive in that disk, and zero outside it. The decomposition is arranged so
that each phase has a positive transverse Hessian for the metric immediately
preceding its addition.

Second, let $F$ be the current immersion and consider one addition
$a^2dx^2$. Write $y$ for a complementary coordinate. The relevant plane is
perpendicular to both $F_y$ and $F_{yy}$.
A periodic path on a circle in this plane has the prescribed average and
adds $a^2dx^2$ to the leading metric. Orthogonality to $F_{yy}$ removes
the first transverse metric error. A finite expansion then cancels the
oscillatory errors to an arbitrarily prescribed order. Its remaining mean
terms form a finite differential perturbation of the induced-metric map.
Small corrections on a slower scale reduce this perturbed map's error to
the required finite accuracy.

Third, turn the resulting approximate metric into an exact one. We prove a
smooth correction theorem under the pointwise hypothesis that the
vector-valued second fundamental form is nonzero. A finite construction solves
the linearized equation on individual oscillatory modes to arbitrarily high
finite accuracy, with suitably chosen phase directions.
Smoothing between successive steps controls the lost derivatives; increasing
the accuracy in successive blocks gives convergence in every smooth norm.
No global inverse of the linearized metric map is required.

Finally, retain the geometric conditions needed for the remaining disks.
The large oscillation may rotate the second fundamental form. Conditions on
the later boundary curves, and a choice of large transverse derivatives at the
turns of the velocity loop, prevent its relevant normal directions from
vanishing or becoming opposite to the preferred normal. The small $C^2$ exact correction preserves the
required derivative asymptotics. Together with the exact Gauss equation,
these asymptotics establish the strict boundary and crossing conditions.
After the finite list of additions, the induced metric is $g$ by
\eqref{eq:intro-sum}.

The two analytical devices used here have separate roles. The small-mode
construction supplies high-order metric accuracy while its displacement is
small in $C^2$. The large circle construction changes the metric by a prescribed
positive primitive tensor while retaining selected second-derivative
information. Their combination is what permits a finite global construction
in codimension two.  The number of global primitive additions is finite;
within each addition, however, the exact correction uses an infinite
smoothing iteration converging to one map in every $C^m$ norm.

\subsection{Organization}
Section~\ref{sec:preliminaries} sets up normal geometry and admissible phases.
Section~\ref{sec:primitive-statement} states the primitive-addition theorem,
including the boundary conditions it preserves.
Sections~\ref{sec:small-modes} and~\ref{sec:exact-correction} prove the small
increment and exact correction results.
Sections~\ref{sec:velocity-loop} and~\ref{sec:primitive-realization} construct
the velocity loop and realize the primitive metric.
Section~\ref{sec:geometry-preservation} completes the primitive-addition proof.
Section~\ref{sec:global-construction} prepares the finite list and proves
Theorem~\ref{thm:main}.
